\documentclass[10pt,a4paper]{amsart}
\usepackage[margin=19mm]{geometry}
\usepackage{amsmath,amssymb,mathtools}
\usepackage[scr=rsfs,cal=euler]{mathalfa}
\usepackage{xfrac}
\usepackage[unicode,hidelinks,bookmarks=false]{hyperref}
\newcommand{\ie}{i.e.\ }
\newcommand{\cf}{cf.\ }
\newcommand{\resp}{respectively\ }
\newcommand{\Subsection}{\subsection}
\providecommand{\nobreakdash}{\nobreak\hbox{-}}
\setlength{\emergencystretch}{2em}
\multlinegap0pt
\usepackage{amssymb}
\usepackage{mathtools}
\usepackage{xfrac}
\usepackage[shortlabels]{enumitem}
\usepackage{xcolor}
\newtheorem{prop}{Proposition}
\newtheorem{lemma}{Lemma}
\title{The least quadratic residue and \\ integers represented by quadratic forms}
\author{Kannan Soundararajan, Jo\~ao C. C. Vargas}
\date{Source: arXiv:2607.29566v1 (CC BY 4.0)}
\begin{document}
\maketitle

\noindent\emph{Source and licence.} The least quadratic residue and \\ integers represented by quadratic forms. Version arXiv:2607.29566v1 (CC BY 4.0), distributed by arXiv under the Creative Commons Attribution 4.0 International licence, http://creativecommons.org/licenses/by/4.0/, as recorded in the arXiv licence metadata for this exact version and shown on the abstract page https://arxiv.org/abs/2607.29566v1. Full licence notice: \texttt{licenses/arxiv-2607.29566-CC-BY-4.0.txt}.\par\smallskip

\noindent\textit{Editorial scope.} Counted unit: the article's prop prop2.1 (source0.tex lines 238--243). The article's complete original proof of it is delivered with it (source0.tex lines 280--307, one \texttt{proof} environment of 28 lines, which the article places away from the statement). No mathematics is written, repaired or re-derived: the statement and the proof are the article's own lines, in the article's own notation, and no retained line is substituted or re-flowed.  Retained uncounted as the source-local context the proof needs: lines 250--256.  Nothing was omitted from any retained span, so the package carries no omission record.  No retained line contains a citation, so the wrapper carries no bibliography and no bibliography entry.  No retained line contains a figure environment, an image-inclusion command or drawing code, so no image is delivered and the package declares no asset dependency.  Editorial formatting only, in addition: an AMS \texttt{amsart} preamble that restates the article's own declarations and macros used by the retained lines, namely the environments \texttt{prop}, \texttt{lemma} on separate counters, together with the standard packages those lines need.  The retained environments print in this document as Proposition 1, Lemma 1 in order of appearance, whereas the article prints the same items with the numbering of its own full document.  The complete native source of the article as distributed in the arXiv e-print for this exact version is bundled unchanged as \texttt{sources/206499/source0.tex}.
\begin{lemma}\label{newtonsIdent} 
    Let $k$ be a natural number and let ${\mathbb F}_2$ denote the field with two elements, and $\overline{\mathbb F}_2$ its algebraic closure.  Suppose $z_1$, $\ldots$, $z_n$ are distinct non-zero elements of $\overline{{\mathbb F}}_2$ with 
    $$ 
    z_1^{2j-1} + z_2^{2j-1} + \ldots + z_n^{2j-1} =0
    $$
    for all $1\le j\le k$.   Then we must have $n> 2k$.  
\end{lemma}

\begin{prop} \label{prop2.1}  For each integer $k\ge 2$ there exist completely multiplicative functions $f_1$, $\ldots$, $f_k : {\mathbb N} \to \{-1,1\}$ such that the least square-free integer $r>1$ with $f_j(r)=1$ for all $1\le j\le k$ satisfies 
$$ 
r \ge 4^{k-Ck/\log k}, 
$$
for a suitable absolute constant $C$. 
\end{prop} 

\begin{proof}[Proof of Proposition \ref{prop2.1}]   We first select a bijection between the primes and irreducible polynomials in ${\mathbb F}_2[x]$ apart from the polynomial $x$ as follows.  Running over the primes in ascending order, associate to $p$ an irreducible polynomial of smallest degree in ${\mathbb F}_2[x]$ that has not already been used (and omitting $x$).  For instance, associate $2$ to the polynomial $x+1$, $3$ to $x^2+x+1$, $5$, $7$ and $11$ to the three irreducible polynomials of degree $3$ (in any way), and so on.  Let $Q_p$ denote the irreducible polynomial associated to $p$.  Since $\pi(x) \sim x/\log x$, and the number of irreducibles in ${\mathbb F}_2[x]$ of degree $d$ is $\sim 2^d/d$, we see that 
$$ 
\text{deg}(Q_p) = \frac{\log p}{\log 2} + O(1). 
$$
Let ${\mathcal Z}_p$ denote the set of roots of $Q_p$ in $\overline{\mathbb F}_2$, and note that, since the irreducible $x$ has been omitted, the different sets ${\mathcal Z}_p$ are disjoint subsets of non-zero elements of $\overline{\mathbb F}_2$.  

For each natural number $j$, and prime number $p$, define 
$$
e_j(p) = \sum_{z\in {\mathcal Z}_p} z^{2j-1},  
$$
which is an element of ${\mathbb F}_2$.   Define a completely multiplicative function $f_j$ by setting on primes $p$
$$ 
f_j(p) = (-1)^{e_j(p)}.  
$$
We claim that the proposition holds for this choice of completely multiplicative functions.  Suppose that $r\ge 2$ is a square-free number with $f_j(r)=1$ for all $1\le j\le k$.  Put ${\mathcal Z} =\cup_{p|r} {\mathcal Z}_p$, and note that by construction 
$$
\sum_{z\in \mathcal{Z}} z^{2j-1} = \sum_{p|r} \sum_{z\in \mathcal{Z}_p} z^{2j-1} = \sum_{p|r} e_j(p) = 0, 
$$
for all $1\le j\le k$.   Therefore, by Lemma \ref{newtonsIdent} we obtain 
$$ 
2k < |{\mathcal Z}| = \sum_{p|r} \text{deg} (Q_p) = \sum_{p|r} \Big(\frac{\log p}{\log 2} + O(1)\Big) = \frac{\log r}{\log 2} + O(\omega(r)),  
$$
where $\omega(r)$ denotes the number of distinct prime factors of $r$.  Since $\omega(r) =O(\log r/\log \log r)$, it follows that 
$$ 
\log r \ge k\log 4 - Ck/\log k,
$$
for a suitable absolute constant $C$, which completes our proof. 
\end{proof}

\end{document}
